% Standalone problem document, generated from the adjacent README.md.
% Compile with XeLaTeX. Mathematical content is maintained in Markdown.
\documentclass[11pt,a4paper]{article}
\usepackage[fontsize=10.5pt]{scrextend}
\usepackage[margin=22mm,headheight=15pt,headsep=9mm,footskip=11mm]{geometry}
\usepackage{fontspec}
\setmainfont{texgyrepagella}[Extension=.otf,UprightFont=*-regular,BoldFont=*-bold,ItalicFont=*-italic,BoldItalicFont=*-bolditalic]
\setsansfont{texgyreheros}[Extension=.otf,UprightFont=*-regular,BoldFont=*-bold,ItalicFont=*-italic,BoldItalicFont=*-bolditalic]
\setmonofont{texgyrecursor}[Extension=.otf,UprightFont=*-regular,BoldFont=*-bold,ItalicFont=*-italic,BoldItalicFont=*-bolditalic,Scale=MatchLowercase]
\usepackage{amsmath,amssymb}
\usepackage{unicode-math}
\setmathfont{texgyrepagella-math.otf}
\usepackage{xcolor,microtype,enumitem,needspace,fancyhdr}
\usepackage{bookmark}
\definecolor{ink}{HTML}{17344B}
\definecolor{accent}{HTML}{197D80}
\definecolor{muted}{HTML}{596773}
\hypersetup{colorlinks=true,linkcolor=accent,urlcolor=accent,pdftitle={SP-09 - Unitary-orbit
distance under finite block
repetition},pdfauthor={Open Problems in Numerical Linear Algebra}}
\urlstyle{same}
\setlength{\parindent}{0pt}
\setlength{\parskip}{5pt plus 1pt minus 1pt}
\linespread{1.04}
\setlength{\emergencystretch}{3em}
\widowpenalty=10000
\clubpenalty=10000
\displaywidowpenalty=10000
\allowdisplaybreaks[1]
\setlist{nosep,leftmargin=1.5em}
\providecommand{\tightlist}{\setlength{\itemsep}{0pt}\setlength{\parskip}{0pt}}
\providecommand{\pandocbounded}[1]{#1}
\pagestyle{fancy}
\fancyhf{}
\fancyhead[L]{\sffamily\footnotesize\color{muted}OPEN PROBLEMS IN NUMERICAL LINEAR ALGEBRA}
\fancyhead[R]{\sffamily\bfseries\footnotesize\color{accent}SP-09}
\fancyfoot[L]{\parbox[b]{0.9\textwidth}{\sffamily\fontsize{8}{10}\selectfont\color{muted}Editorial ratings; a bounded search cannot certify that a problem remains unsolved.\\Literature check: 2026-09-10}}
\fancyfoot[R]{\sffamily\footnotesize\color{muted}\thepage}
\renewcommand{\headrulewidth}{0.3pt}
\renewcommand{\headrule}{\hbox to\headwidth{\color{accent}\leaders\hrule height \headrulewidth\hfill}}
\makeatletter
\renewcommand{\section}{\Needspace{5\baselineskip}\@startsection{section}{1}{0pt}{12pt plus 2pt minus 2pt}{4pt}{\sffamily\bfseries\large\color{ink}}}
\renewcommand{\subsection}{\Needspace{4\baselineskip}\@startsection{subsection}{2}{0pt}{9pt plus 1pt minus 1pt}{3pt}{\sffamily\bfseries\normalsize\color{ink}}}
\makeatother
\setcounter{secnumdepth}{0}
\begin{document}
{\sffamily\small\color{muted}Eigenvalues and inverse problems\par}
\vspace{3pt}
{\raggedright\hyphenpenalty=10000\exhyphenpenalty=10000\sffamily\bfseries\fontsize{21}{24}\selectfont\color{ink}Unitary-orbit
distance under finite block repetition\par}
\vspace{7pt}
{\sffamily\small\textbf{Difficulty:} challenging\quad\textbf{Importance:} interesting
to specialist\par}
{\sffamily\small\bfseries\color{accent}Status: Partially resolved\par}
\vspace{4pt}
{\color{accent}\hrule height 0.8pt}
\vspace{6pt}
\textbf{Rating rationale:} Challenging reflects preservation of a
nontrivial spectral-norm orbit distance under larger unitary mixing;
specialist impact concerns finite matrix amplification and
operator-algebraic nearness.

\subsection{Problem statement}\label{problem-statement}

For \(A,B\in\mathbb C^{n\times n}\), define the distance between their
unitary similarity orbits in the spectral norm by

\[
\delta_n(A,B)=\min_{U\in\mathbb C^{n\times n},\ U^*U=I_n}
\|A-U^*BU\|_2.
\]

For a positive integer \(k\), write \(A^{(k)}=I_k\otimes A\), the block
diagonal matrix with \(k\) copies of \(A\). Is it true that, for every
\(n\ge3\), every finite integer \(k\ge2\), and every pair of normal
matrices \(A,B\in\mathbb C^{n\times n}\),

\[
\delta_{nk}\bigl(A^{(k)},B^{(k)}\bigr)=\delta_n(A,B)?
\]

Normality means \(AA^*=A^*A\) and \(BB^*=B^*B\). The minimizing unitary
on the left is unrestricted in dimension \(nk\); it may mix different
copies. Thus the question asks whether block repetition can improve the
best unitary-similarity fit between normal matrices.

\subsection{Why it matters}\label{why-it-matters}

This is a matrix-nearness problem with prescribed spectra. It tests
whether a structured enlargement changes the optimum in spectral norm, a
metric for which matching eigenvalues alone need not give the orbit
distance.

\newpage
\subsection{References}\label{references}

L. W. Marcoux and Y. Zhang,
\href{https://doi.org/10.1016/j.jfa.2020.108778}{\emph{On Specht's
Theorem in UHF C}-algebras*}, Journal of Functional Analysis 280 (2021),
108778, §5, Question 5.4, pp.~26--27;
\href{https://uwspace.uwaterloo.ca/items/5359b050-4379-49b2-bdda-66ac3611d3ca}{institutional
full text}.

L. W. Marcoux, P. Sarkowicz, and Y. Zhang,
\href{https://arxiv.org/html/2508.13834v1}{\emph{Kaplansky's problem and
unitary orbits in matrix amplifications}}, arXiv:2508.13834v1 (2025),
§3.5 and Propositions 3.15, 3.21, and 4.14.

\subsection{Earlier status check ---
2026-09-08}\label{earlier-status-check-2026-09-08}

Theorem 5.3 of the source proves equality for normal matrices of order
two. Theorem 5.1 establishes a strict decrease for some unrestricted
matrices, so normality cannot simply be omitted. The example after
Question 5.4 uses infinitely many copies and does not resolve the
finite-\(k\) question. The direct follow-up 2508.13834v1 records the
order-two result after Corollary 3.5. The corollary itself proves
equality for self-adjoint matrices of every order. Its normal-element
Propositions 3.15 and 3.21 concern zero orbit distance; its Proposition
4.14 counterexamples have no normality requirement. Neither settles
preservation of all distances for finite normal matrices. The published
§5 and these later passages were checked. Searches for
\textit{Marcoux Zhang Question 5.4 normal matrices},
\textit{normal matrices ampliations unitary orbit distance}, and
2025/2026 follow-ups found no finite-dimensional normal counterexample
or general proof. This is limited search evidence, not a certificate of
openness.

\subsection{Audit update --- 2026-09-10}\label{audit-update-2026-09-10}

Rechecked
\href{https://arxiv.org/html/2508.13834v1}{Marcoux--Sarkowicz--Zhang},
Corollary 3.5: equality is proved for every pair of self-adjoint
matrices, in every dimension and finite multiplicity. This substantive
normal subclass was missing from the earlier status explanation. General
normal matrices remain unresolved in the checked sources; finite-normal
amplification searches found no completion. The nonnormal
counterexamples and zero-distance results do not settle the remaining
positive-distance question.
\end{document}
